Insect pollinators contribute to the production of roughly three quarters of the world's food crops~\cite{klein2007pollinators}, and their populations are under sustained pressure from habitat loss, agrochemical exposure, and climate change~\cite{potts2010decline}. Greenhouse operators in particular have to compensate for the under-supply of pollinators in enclosed environments, and several research groups have proposed autonomous robotic pollinators as a partial answer~\cite{strader2019pollination}. The perception system on such a robot has one job that recurs every time the operator wants to extend it to a new crop: detect every open flower in a row of plants reliably, in real time, on hardware that fits on the robot.

Detecting flowers across a single, fixed crop is by now a routine deep-learning task~\cite{dias2018apple, tian2019apple, genemola2020fruit}. Detecting them across an arbitrary new crop, with only a handful of labels for that crop, is much harder. Fully retraining a detector from scratch needs hundreds of labeled images and tends to fail badly with fewer than fifty. Fine-tuning from a generic checkpoint such as the COCO-pretrained YOLOv8~\cite{lin2014coco, jocher2023ultralytics} works better because the convolutional features transfer, but the model still has to learn the concept of "flower" from scratch on top of features tuned for cars, dogs, and indoor furniture. We hypothesized that a checkpoint pretrained on a small but \emph{flower-specific} cross-crop dataset would transfer better at small budgets than COCO.

This paper tests that hypothesis empirically. We construct a joint multi-crop flower dataset by harmonizing the annotations of three open Roboflow datasets (apple, strawberry, tomato) into a single \texttt{flower} class, pretrain YOLOv8n on the joint set, and then fine-tune on a held-out target crop (kiwi) at five label budgets between 10 and 181 with three random seeds each. We compare this protocol --- which we call CrossPoll --- against three baselines: training from scratch, fine-tuning from COCO, and fine-tuning from the joint pretrain without backbone freezing. The full sweep comprises 72 fine-tune runs.

The story that emerges is more nuanced than we expected. Fine-tuning the joint pretrain without freezing the backbone produces a sharp accuracy collapse at $n{=}25$: the validation mAP peaks at epoch 3, deteriorates monotonically for the next fifteen epochs, and early-stops near zero. Per-epoch logs make the failure mode visible --- the joint pretrain's flower-specific features are fragile under the noisy gradient updates that 25 images and a batch size of 16 produce. Freezing the YOLOv8n backbone (modules 0--9 in the Ultralytics convention) protects those features and rescues the small-budget regime. With backbone freezing, CrossPoll attains 0.630 test mAP@0.5 at $n{=}10$ versus 0.362 for COCO transfer, and 0.871 at $n{=}25$ versus 0.892 for COCO. From $n{=}50$ upward, COCO transfer regains a small lead because the backbone breadth advantage of full fine-tuning starts to pay off.

We make four contributions:

\begin{enumerate}
\item We harmonize three open agricultural datasets into a single-class joint flower training corpus and release the harmonization script (\texttt{scripts/harmonize\_labels.py}) so that the dataset can be reproduced from the public sources.
\item We document a small-budget failure mode of CrossPoll-style fine-tuning that does not appear in standard YOLOv8 transfer-learning recipes, and we trace it to small-batch gradient noise on a domain-specialized backbone.
\item We show that backbone freezing (\texttt{freeze=10}) eliminates the failure mode and produces the strongest results we have at $n \leq 25$ labels.
\item We present a clean budget-by-method comparison across four methods, five budgets, and three seeds, with all numbers grounded in a single reproducible \texttt{summary.csv} and the table and figure regenerated from that file by a single script.
\end{enumerate}

The paper is structured as follows. Section~\ref{sec:related} situates CrossPoll relative to transfer learning, few-shot detection, and agricultural computer vision. Section~\ref{sec:method} describes the harmonization, joint pretraining, and fine-tuning protocols. Section~\ref{sec:experiments} reports the experimental setup. Section~\ref{sec:results} presents the headline numbers. Section~\ref{sec:discussion} interprets them and lists the limitations honestly. Section~\ref{sec:conclusion} closes with the next experiments.
